% !TITLE 少年行四首
% !AUTHOR 王维
% !DYNASTY 唐
% !GENRE 七言绝句
% !ORDER 36

\begin{poem}
新丰美酒斗十千\textsuperscript{1}，咸阳游侠多少年。
相逢意气为君饮\textsuperscript{2}，系马高楼垂柳边。

出身仕汉羽林郎\textsuperscript{3}，初随骠骑\textsuperscript{4}战渔阳\textsuperscript{5}。
孰知不向边庭苦\textsuperscript{6}，纵死犹闻侠骨香\textsuperscript{7}。

一身能擘两雕弧\textsuperscript{8}，虏骑千重只似无。
偏坐金鞍调白羽\textsuperscript{9}，纷纷射杀五单于\textsuperscript{10}。

汉家君臣欢宴终，高议云台\textsuperscript{11}论战功。
天子临轩赐侯印\textsuperscript{12}，将军佩出明光宫\textsuperscript{13}。
\end{poem}

\begin{annotations}
\notetext{1}{新丰美酒斗十千：新丰（长安附近地名）的美酒一斗值十千钱。斗，酒器，也是计量单位。}
\notetext{2}{意气为君饮：意气相投，就为你喝这杯酒。不需要交情深浅，意气相合就已足够。}
\notetext{3}{羽林郎：羽林军（汉代禁卫军）的军官。王维以汉比唐。}
\notetext{4}{骠骑：指汉代骠骑将军霍去病，少年名将。此处代指唐军将领。}
\notetext{5}{渔阳：今北京一带，唐代安禄山起兵之处，泛指边塞战场。}
\notetext{6}{孰知不向边庭苦：谁不知道边塞艰苦？但少年游侠不因苦而退缩。孰知，谁不知道。}
\notetext{7}{侠骨香：侠客的骨头是香的——即便战死，侠名也会流传千古。}
\notetext{8}{擘两雕弧：拉开两张雕刻花纹的弓。擘（bò），拉开；雕弧，雕有花纹的硬弓。}
\notetext{9}{偏坐金鞍调白羽：侧身坐在金马鞍上，调整白羽箭。偏坐，侧身坐（射箭姿势）。白羽，箭名。}
\notetext{10}{五单于：五个单于（匈奴首领），这里泛指敌军将领。单于（chán yú）。}
\notetext{11}{云台：汉代洛阳宫中的高台，陈列功臣画像。}
\notetext{12}{临轩赐侯印：天子亲临殿前，赐予侯爵的金印。}
\notetext{13}{明光宫：汉代宫殿名，代指唐朝皇宫。}
\end{annotations}

\begin{appreciation}
这四首诗是王维早年作品，写的是一个少年游侠从长安饮酒到边塞立功、最后受封赏的完整故事。四首连读，像一部短篇小说的四个章节。

第一首：新丰美酒斗十千，咸阳游侠多少年——新丰的美酒一斗值十千钱，咸阳城里的游侠少年不知有多少。新丰是长安附近产美酒的地方，斗十千形容酒贵，也衬托出少年们的豪奢。相逢意气为君饮，系马高楼垂柳边——意气相投便为你痛饮，把马系在高楼边的垂柳下。这是少年游侠的日常：酒、马、高楼、垂柳，还有随时可以结交的朋友。意气二字，是少年人的通行证，不需要身份，不需要理由，一杯酒就是朋友。

第二首：出身仕汉羽林郎，初随骠骑战渔阳——出身做了汉朝的羽林郎，初次跟随骠骑将军出征渔阳。羽林郎是皇帝的禁卫军，骠骑是汉代名将霍去病的官职。王维用汉事写唐事，是古诗常用的手法。孰知不向边庭苦，纵死犹闻侠骨香——谁知道不去边塞受苦呢？纵然战死，侠骨依然留香。擘（bò）是拉开的意思，雕弧（diāo hú）是雕刻有花纹的弓。不向边庭苦，是不怕边塞之苦；侠骨香，是少年人对荣誉的信仰。

第三首：一身能擘两雕弧，虏骑千重只似无——一个人能拉开两张弓，面对千重敌骑却像没有一样。这是武艺的巅峰，也是勇气的极致。偏坐金鞍调白羽，纷纷射杀五单于——侧身坐在金鞍上，调试白羽箭，纷纷射杀敌军首领。白羽是箭名，五单于是匈奴首领的称号，这里泛指敌军将领。这一首是战斗的高潮，少年从游侠变成了英雄。

第四首：汉家君臣欢宴终，高议云台论战功——汉家君臣欢宴结束，在云台高阁上论功行赏。云台是汉代陈列功臣画像的地方。天子临轩赐侯印，将军佩出明光宫——天子亲临殿前，赐予侯爵印信；将军佩戴着印信，走出明光宫。这是故事的结局：少年功成名就，受封侯爵。

但四首连读，你会发现一种隐隐的失落。第一首的垂柳、美酒、意气，是最动人的；到了第四首，侯印、明光宫，反而显得空洞。王维或许也在暗示：少年人的侠骨香，终究要变成成年人的功名簿。那个系马高楼垂柳边的少年，和那个佩印走出明光宫的将军，还是同一个人吗？
\end{appreciation}
